\chapter{The HHS problem}\label{ch:hhs}

This chapter fixes the object of study for the rest of Part~\ref{part:results}. Hook, Huang and Shalaby (HHS)~\cite{HHS2026} argue that the resonant conversion of dark-photon dark matter into plasma oscillations saturates at a level of the order of the electron thermal energy, and that the cosmological bounds built on the opposite assumption therefore fail. The argument rests on one-dimensional particle-in-cell calculations, and the present note exists because we wanted to know which parts of it a noise-free continuum solver can confirm, and which parts it cannot. Before any of our own results are discussed we need to be exact about what HHS compute, which of their curves we compare with, how the energy they plot relates to the energies our solvers carry, and how a number we first took from their figure turned out to be wrong by a factor of two. The last point is told without softening, because several later chapters hinge on it.

\section{Set-up and claims of HHS}

The physical premise is the one used in all earlier work on resonant conversion. A dark photon of mass $m_{A'}$ with kinetic mixing $\epsilon$ behaves, for the plasma, like a nearly uniform oscillating electric field, because dark matter is cold and the ratio $k/\omega$ of the field is of order $H/\omega_p\ll1$ at the epoch of interest. When $\omega_p=m_{A'}$ the field drives the electrons coherently, and the conversion probability of the linear theory is the Landau--Zener expression of their Eq.~(2), $P_{A'\to\gamma}=\pi\epsilon^2 m_{A'}^2\,[\mathrm d\log\omega_p^2/\mathrm d\chi]^{-1}_{\omega_p=m_{A'}}$, with $\chi$ the line-of-sight distance. The HHS paper does not dispute this formula. It disputes the assumption that the linear response persists until the resonance has been crossed.

Their criterion for the breakdown of linear theory is given in Eq.~(3). The driven electrons carry a current $j\simeq\omega_pE$ per unit charge density, and the perturbative expansion fails when the collective drift exceeds the thermal speed, that is when $eE\simeq m_e\omega_pv_{th}^e$. They define the quiver velocity $v_q=eE/(m_e\omega_p)$ and call $v_q/v_{th}^e$ the measure of nonlinearity. At resonance the accumulated quiver velocity grows linearly, $v_q\simeq v_q^D\omega_pt/2$ (their SM Appendix~A), where $v_q^D=\epsilon eE_0/(\omega_pm_e)$ is the instantaneous strength of the dark-photon drive of their Eq.~(5). The quantity we call $\vq$ throughout this note is $v_q^D/v_{th}^e$. Because it is the \emph{instantaneous} drive, and the accumulated amplitude grows with time, the same $\vq$ reaches the threshold $v_q=v_{th}^e$ at $\omega_pt=2v_{th}^e/v_q^D$, so that $\vq=0.03$ gets there at $\omega_pt\simeq67$ and $10^{-3}$ at $\omega_pt\simeq 2000$.

Beyond that point HHS describe a chain of events that they support with the Zakharov equations and with simulation. The ponderomotive potential of their Eq.~(4), $\Phi_p=e^2\hat E^2/(4m_e\omega^2)$, pushes both species towards regions of low field amplitude. The $k\simeq0$ Langmuir wave is unstable to the sub- and supersonic modulational instability and to the electrostatic decay instability, which populate Langmuir and ion-acoustic waves over a range of $k$ from the ion plasma scale up to the inverse Debye length. The ion-acoustic perturbations modulate the total density on a longer time scale, hence the plasma frequency, and thereby detune the uniform resonance. For $\vq\gtrsim0.01$ they distinguish a fast regime in which the electron dynamics becomes non-perturbative almost at once and the field energy overshoots the thermal energy before the ions respond; for $\vq\lesssim0.01\simeq(m_e/m_i)^{1/2}/2$ a slow regime in which the nonlinearity sets in before the Langmuir energy reaches the thermal energy. In both regimes the final state is said to be a mixture of Langmuir and ion-acoustic waves with electrons about thirty times hotter than the ions, and the transfer is said to stop. The conclusion in the abstract is that the deposited energy is limited to about the thermal energy of the electrons at the time of conversion, and that the bounds weaken by factors from $3000$ to $10^{7}$.

The simulations that carry this conclusion are described in the Supplemental Material (SM). The code is SHARP, adapted to a one-dimensional electron--ion plasma with $m_i/m_e=1836$; the pusher is the Vay algorithm, the deposition and interpolation use fifth-order splines, and the box is $L=40\,c/\omega_p$ with $N_x=1000$ cells of $0.04\,c/\omega_p$. Both species start from Maxwellians with $v_{th}^e=3.1\times10^{-2}c$, that is $k_BT=10^{-3}m_ec^2$, uniformly distributed in space, and the drive is a prescribed uniform field $A_0\cos\omega t$ with no back-reaction on the dark photon. The number of particles per cell ranges from $2\times10^2$ to $2\times10^5$, selected with the shot-noise estimate of their Eq.~(S37) so that the noise energy sits below the thermal energy but the driven energy overtakes it at $\omega_pt\simeq2.4$ (their Eq.~S38). Two resonant cases are shown in their Fig.~2, the upper panel with $\vq=0.03$ and the lower with $10^{-3}$, each filtered with a moving average over $200\,\omega_p^{-1}$. Landau--Zener passages are in their Fig.~3, which we do not address here beyond the discussion of inhomogeneous conversion in Chapter~\ref{ch:nonrel}. The summary figure of the revised bounds is their Fig.~4.

Two features of the Fig.~2 narrative will matter later. The first is that the black curve, the energy in $k\neq0$ modes and electron heating, tracks where the linear-theory curve would have been: ``at late time, the black line matches where the red line was'', in the sense that an order-unity fraction of the energy has moved from the $k=0$ wave to higher-$k$ motion and thermal energy. The second is the position of the departure itself, which for $\vq=0.03$ is at several hundred inverse plasma frequencies. Both are statements about a simulation with a given noise level, a given shape function and a given pusher, and we shall see that each of those choices needs to be controlled before the statement can be attributed to a mechanism.

\section{Energies: theirs and ours}

HHS plot, in Figs.~2 and~3, the quantity they call $\Delta E$, in units of $n_em_ec^2$. According to the SM (Eqs.~S35 and~S36 and the discussion after them) this is the change of the thermal energy density of each species,
\begin{equation}
E_{\rm th}^{\alpha}=\frac{n_\alpha m_\alpha (v_{th}^\alpha)^2}{2}=\frac{n_\alpha k_BT_\alpha}{2},\qquad
(v_{th}^\alpha)^2=\int\mathrm dv_x\,v_x^2f_\alpha(v_x),
\label{eq:hhs-eth}
\end{equation}
where the average velocity in $x$ is zero at $t=0$ and the integral is over the one-dimensional velocity distribution. The factor $\tfrac12$ matters: with $k_BT=10^{-3}m_ec^2$ the initial value is $5\times10^{-4}n_em_ec^2$, which is what the dotted gray line of their Fig.~2 marks and what we call $\Eth$. A footnote of the SM adds that evolving all three velocity components would sum the three directional temperatures; the calculation is one-dimensional in velocity, so $\Eth$ is the energy of one degree of freedom. Our reading of the plotted curves, which we take from the main text and the SM together, is that the thermal energy is computed after removing the $k=0$ component, so that a coherent drift of the whole electron population (the quiver motion) is excluded, while a coherent $k\neq0$ wave is included. This is an inference, not a quotation, and it is part of the open question of Chapter~\ref{ch:discussion} whether the plotted quantity contains anything else.

The energies carried by our own solvers are those of the work ledger of the paper~\cite{DarkSPECTRAX}. Writing $K_e$ for the electron kinetic energy, $\Wext=\int\bar J\,E_{\rm drive}\,\mathrm dt$ for the work done by the prescribed drive on the plasma, and $\bar E$ for the $k=0$ electric field, the energy theorem reads
\begin{equation}
K-K_0=W_\gamma+W_D+\Wext,\qquad U_\gamma-U_{\gamma0}=-W_\gamma,
\label{eq:ledger}
\end{equation}
and in the prescribed-drive mode of this note $W_D$ is absent. The ledger holds to solver tolerance in the nonrelativistic Hermite code and, with the rest-mass term $(\gamma-1)m c^2$ in $K$, in the relativistic one. Our records measure works in units of $n T_e V$, with $V$ the box volume; since $\Eth=\tfrac12nT_e$ for one degree of freedom, one unit of $nT_eV$ is two thermal energies. For example the window-end work of $381\,nT_eV$ at $\vq=0.07$ (Chapter~\ref{ch:nonrel}) corresponds to $762\,\Eth$ in HHS units, almost all of it coherent quiver motion. The consequence is that the headline numbers of the continuum studies (works of tens to hundreds of $nT_eV$) are not comparable to the HHS ordinate, which excludes the coherent motion, and that a separate diagnostic is needed.

That diagnostic is $\Delta E_e^{1D}$: the change, relative to $t=0$, of $\tfrac12 n m_e$ times the one-dimensional electron velocity variance, evaluated about the box-mean flow, so that the $k=0$ drift drops out, and then averaged over the trailing $200\,\omega_{pe}^{-1}$ in the same way HHS filter their curves. In the Hermite code it follows from the zeroth to second Hermite coefficients, and the box mean is the $k=0$ component; in the grid code it is a velocity integral of $f$; in the PIC it is the particle variance minus the cell-averaged flow. All curves in Chapters~\ref{ch:nonrel} and~\ref{ch:rel} that are compared with HHS use this measure and the unit $\Eth=5\times10^{-4}n_em_ec^2$. The other two energies that HHS show, the mean field $\langle E^2\rangle/2$ and the ion thermal energy, have direct counterparts: the first is $\bar E^2/2$ in our records, the second the ion variance.

A comparison between the mean-field energy in the two frameworks needs one more remark. The mean field $\bar E$ is the response of the plasma to the driver, not the driver itself, and the exact linear law for it, $\ddot{\bar E}+\omega_L^2\bar E=-\omega_L^2E_0\cos\omega t$ with $\omega_L^2=\omega_{pe}^2(1+m_e/m_i)$, is the baseline against which the departure is measured. In terms of it the energy delivered to the plasma is the linear resonant law $\Wlin(t)$. The statement ``the resonance is lost at time $t_c$'' is then operationalised in our work as the first time at which $\Wext(t)$ differs from $\Wlin(t)$ by $10\,\%$, with the further requirement that the departure be confirmed in refinement. The corresponding HHS statement is a visual one, made on the mean-field curve of the same figure.

\section{The re-digitization of Fig.~2, and the earlier misreading}

All of the quantitative comparison with HHS rests on one number, the value of the black curve of the upper panel of Fig.~2 at $\omega_pt=1000$. We obtained it twice, and the two readings differ by a factor of two.

The first reading was made early in the relativistic PIC study, before the convergence ladder had been run. It was a read-off of the black curve at $\omega_pt=1000$ from the logarithmic ordinate, and it gave approximately $25\,\Eth$; the procedure was not recorded in enough detail to be repeated. The relativistic PIC of Chapter~\ref{ch:rel} then gave $21.5\,\Eth$ at 103 particles per cell, within $15\,\%$, and the first status entries of the repository say that the relativistic run is ``within 8--15\,\% below the digitized HHS Fig.~2 curve''. That agreement was taken as evidence that the relativistic PIC reproduces HHS and that the Newtonian pusher, which gave $56\,\Eth$, overestimates them by a factor of $2.2$. The first PIC study also declared convergence gates, and two of them failed: the heating changed by $+19\,\%$ when the particles per cell went from 103 to 412 and by about the same when the cell was halved. A follow-up study declared a denser ladder in particle number, cell size, shape function and time step, and that ladder converged. Only after it had converged, and the plateau sat at about half of the number it was being compared with, was the comparison number itself re-examined.

The second reading was made from the vector figure of the article, rendered at 600~dpi, with the axes calibrated on the detected grid lines of the plot rather than on printed tick labels. The dotted line marking the initial thermal energy reads $5.02\times10^{-4}$ against the nominal $5\times10^{-4}$, which checks the calibration to $0.4\,\%$. With this calibration the black solid curve, the $k\neq0$ electron energy change, reads $50.1\,\Eth$ at $\omega_pt=1000$ and $52.7\,\Eth$ at $2000$, with an uncertainty of about $3\,\%$ from the line width (record \texttt{studies/relpic/digitization/hhs\_fig2.json}). The earlier value of $\sim25$ was a misreading by a factor of two. We have not reconstructed how it arose; a logarithmic ordinate invites errors of this kind, but that is a remark and not a finding. The point of recording the error is its consequences, not its cause.

Those consequences are listed in the repository in the status entries of 2026-10-25 and 2026-10-27 and in Table~II of the paper. The statement that the relativistic PIC is within 8--15\,\% of HHS was wrong; it is about a half of the HHS curve, and the declared comparison gate C6 (agreement within $10\,\%$ including the digitization uncertainty) fails. The statement that the Newtonian pusher overestimates the heating by a factor of two is wrong if the comparison is with HHS: at 412 particles per cell it gives $62.7\pm9.4\,\Eth$, which is $1.25$ times the HHS value, and it is the relativistic one that is low. Relativity reduces the heating by a factor of about $2.5$, and that statement does not depend on HHS at all. The statement of the open question changed accordingly. Before the correction the HHS curve and the relativistic PIC agreed, and the remaining puzzle was why the continuum codes did not reproduce it. After the correction the continuum puzzle is unchanged and a second one appears: the relativistic PIC of an independent code, converged in every parameter we varied, lands at half of the published curve.

The correction also changes what can be said about the dependence on ingredients. Newtonian dynamics, which overshoots the relativistic result by $2.5$, brackets HHS from above, and a relativistic code with the HHS shape order and pusher might land anywhere between our plateau and the Newtonian result; this is a matter for the candidate causes enumerated in Chapter~\ref{ch:rel} and not something the digitization can settle. The uncertainty in the number itself is not the issue. A $3\,\%$ line-width uncertainty cannot move $50$ to $25$. The old value is retained in the repository histories (the status entry states ``the earlier $\sim$25 was a misreading'') and in Table~II of the paper, where the first PIC comparison is listed among the negative and superseded results. We repeat it here because a note meant to be read after the paper should record both the correction and the path that led to it, and because a reader who finds a ``$\sim25$'' in one of the earlier reports (\texttt{reports/hhsref.md}, \texttt{reports/relhhs.md}) should know that it refers to the superseded reading.

\section{What is compared in what follows}

For the resonant case $\vq=0.03$ there are five observables, and the later chapters report each of them with the same definitions. The first is the time at which the work $\Wext$ departs from the linear resonant law by $10\,\%$, a stand-in for the HHS statement that resonance is lost near $\omega_pt\simeq800$. The second is the ratio of the mean-field energy to its linear value at $\omega_pt=500$ and $1000$. The third is the HHS measure $\Delta E_e^{1D}/\Eth$ at $1000$ and $2000$. The fourth is the quiver speed implied by the work ledger, $\sqrt{2\Wext/(nm_e)}$ in units of $c$, which decides whether the nonrelativistic model is applicable at all. The fifth is the validity time of each calculation: positivity of the electron distribution ($f_e>-10^{-3}$ in units of its maximum), closure of the energy ledger to $10^{-4}$ of the work, and agreement of refinement pairs within $10\,\%$. The fifth is not an observable of the physics but is what gives the others meaning; Chapter~\ref{ch:discussion} returns to it as a methodological point.

Table~\ref{tab:c8-hhs} summarizes the HHS quantities used as targets and the corresponding quantities in the repository.

\begin{table}[t]
\centering
\caption{HHS targets for the resonant case $\vq=0.03$ and their counterparts in the records. Page and equation labels refer to the PRL~\cite{HHS2026} and its SM.}\label{tab:c8-hhs}
\small
\begin{tabular}{@{}p{4.2cm}p{4.3cm}p{5.2cm}@{}}\toprule
HHS quantity & Where & Counterpart here\\\midrule
$k\neq0$ electron energy change, $50.1\,\Eth$ at $\omega_pt=1000$, $52.7$ at $2000$ & Fig.~2 upper panel, re-digitized & $\Delta E_e^{1D}/\Eth$, trailing $200\,\omega_{pe}^{-1}$ mean\\
Mean-field energy $\langle E^2\rangle/2$, about $10^{-3}$ at $\omega_pt\simeq800$, $\lesssim10^{-4}$ at 1000 (units $n_em_ec^2$) & Fig.~2, read from the log plot & $\bar E^2/2$ over its linear value, and in $\Eth$\\
Resonance lost near $\omega_pt\simeq800$ & Fig.~2 text & 10\,\% departure of $\Wext$ from $\Wlin$\\
Thermal energy $5\times10^{-4}n_em_ec^2$ & Eqs.~S35, S36 & $\Eth=nT_e/2$\\
Particle number $2\times10^2$ to $2\times10^5$ per cell, $\Delta x=0.04$ & Eqs.~S37, S38, SM App.~B & PIC ladder 103--1648 per cell\\
Vay pusher, fifth-order splines & SM App.~B & Boris/Vay-family, quadratic and quintic\\\bottomrule
\end{tabular}
\end{table}
